\section{Regional support of physical bond transformations}
\label{sec:peps_semiregular_parent_support}

The bond transformation of \cite[Section~7]{Schuch2010PEPS} acts on
head--tail endpoint pairs. We describe its regional blocks, including the
one-ended factors on crossing edges, and derive the product support of
regional parent spaces. Throughout the graph statements, $\Gamma$ is a
finite simple graph with ordered vertices. For a region $R$, write $E_R$
for its internal edges, $\partial E_R$ for its crossing edges, and
$E_{R^c}$ for its exterior edges. An edge $f=(t,h)$ has $t<h$.

\subsection{Products, coordinate slices, and multiplicity maps}

\begin{definition}[Mixed products of boundary and internal maps]%
    \label{def:peps_parent_support_mixed_product}
    \lean{TNLean.PEPS.mixedPhysicalProductMatrix,
        TNLean.PEPS.mixedPhysicalProductMap,
        TNLean.PEPS.mixedPhysicalProductMap_apply}
    \leanok
    Let $B,E$ be finite sets, with matrices $F_b:\C^X\to\C^Y$
    and $K_f:\C^P\to\C^Q$. Their mixed product has coefficients
    \begin{align}
        \mathcal M(F,K)_{(\tau,\beta),(\sigma,\alpha)}
        =\prod_{b\in B}(F_b)_{\tau_b,\sigma_b}
        \prod_{f\in E}(K_f)_{\beta_f,\alpha_f}.
        \label{eq:peps_parent_support_mixed_product}
    \end{align}
    When $X,P$ are finite, its action is
    $(\mathcal M(F,K)\psi)(\tau,\beta)
        =\sum_{\sigma,\alpha}\mathcal M(F,K)_{(\tau,\beta),(\sigma,\alpha)}
        \psi(\sigma,\alpha)$.
\end{definition}

\begin{theorem}[Mixed products on coherent sums]%
    \label{thm:peps_parent_support_mixed_action}
    \lean{TNLean.PEPS.mixedPhysicalProductMatrix_conjTranspose,
        TNLean.PEPS.mixedPhysicalProductMap_sum_prod}
    \leanok
    \uses{def:peps_parent_support_mixed_product}
    The adjoint satisfies
    $\mathcal M(F,K)^\dagger=\mathcal M(F^\dagger,K^\dagger)$.
    Assume $X,P$ are finite. For a finite index set $A$, weights $w_a$, and factor vectors
    $x_{b,a}\in\C^X$, $y_{f,a}\in\C^P$, one has
    \begin{align}
        \mathcal M(F,K)
        \left(\sum_{a\in A}w_a
        \bigotimes_{b\in B}x_{b,a}\otimes
        \bigotimes_{f\in E}y_{f,a}\right)
        =\sum_{a\in A}w_a
        \bigotimes_{b\in B}(F_bx_{b,a})\otimes
        \bigotimes_{f\in E}(K_fy_{f,a}).
        \label{eq:peps_parent_support_mixed_action}
    \end{align}
\end{theorem}
\begin{proof}\leanok
    Conjugate each coefficient in
    (\ref{eq:peps_parent_support_mixed_product}) to obtain the adjoint
    formula. For~(\ref{eq:peps_parent_support_mixed_action}), interchange
    the finite sum over $a$ with the input-coordinate sums. For each
    fixed $a$, the independent coordinate sums factor into the products
    defining $F_bx_{b,a}$ and $K_fy_{f,a}$.
\end{proof}

\begin{theorem}[Product ranges from one-coordinate slices]%
    \label{thm:peps_parent_support_product_range}
    \lean{TNLean.PEPS.physicalProductMap_fixed_of_slices_fixed,
        TNLean.PEPS.physicalProductMap_slice_mem_range,
        TNLean.PEPS.mem_range_physicalProductMap_iff}
    \leanok
    Let $S,X,Y$ be finite and $F:\C^X\to\C^Y$ a matrix. For
    $\psi\in\C^{Y^S}$, define
    $\psi_{v,\tau}(s)=\psi(\tau[v\mapsto s])$. Then
    \begin{align}
        \psi\in\operatorname{im}F^{\otimes S}
        \quad\Longleftrightarrow\quad
        \forall v\in S,\ \forall\tau\in Y^S,\quad
        \psi_{v,\tau}\in\operatorname{im}F.
        \label{eq:peps_parent_support_product_range}
    \end{align}
    In the forward implication $Y$ need not be finite.
    For any matrix $P:\C^Y\to\C^Y$ with $Y$ finite, if
    $P\psi_{v,\tau}=\psi_{v,\tau}$ for every $v,\tau$, then
    $P^{\otimes S}\psi=\psi$; idempotence of $P$ is unnecessary.
    The assertions allow empty sets and impose no rank condition on $F$.
\end{theorem}
\begin{proof}\leanok
    A slice of $F^{\otimes S}\chi$ is a linear combination of columns
    of $F$, with coefficients obtained by contracting all other factors.
    For the converse choose a linear right inverse of $F$ on its image
    and extend it to $L:\C^Y\to\C^X$. Thus $FL$ fixes every slice.
    Applying $FL$ at one coordinate fixes $\psi$. Induction over finite
    subsets of $S$ shows that applying it at all coordinates also fixes
    $\psi$. Consequently
    $F^{\otimes S}(L^{\otimes S}\psi)=(FL)^{\otimes S}\psi=\psi$.
    The same induction proves the assertion about an arbitrary $P$.
\end{proof}

\begin{definition}[Endpoint labels and equal-copy amplitudes]%
    \label{def:peps_parent_support_copy_amplitude}
    \lean{TNLean.PEPS.multiplicityEndpointBase,
        TNLean.PEPS.multiplicityEndpointCopy,
        TNLean.PEPS.multiplicityBondAmplitude}
    \leanok
    Let $I$, the sets $X_i$, and the multiplicity sets $M_i$ be finite.
    Put $X=\coprod_iX_i$, $Y=\coprod_i(X_i\times M_i)$, and
    $m_i=|M_i|$. For $y=(i,a,u)\in Y$, define
    $b(y)=(i,a)$ and $c(y)=(i,u)$. Set
    \begin{align}
        a(y,z)=
        \begin{cases}
            m_i^{-1/2}, & y=(i,a,u),\ z=(i,b,u), \\
            0,          & c(y)\ne c(z).
        \end{cases}
        \label{eq:peps_parent_support_copy_amplitude}
    \end{align}
    Whenever the first case occurs, $M_i$ contains $u$, so $m_i>0$.
\end{definition}

\begin{theorem}[Sparse action of multiplicity restoration]%
    \label{thm:peps_parent_support_sparse_action}
    \lean{TNLean.PEPS.endpointEmbeddingMatrix_conjTranspose_mulVec,
        TNLean.PEPS.multiplicityBondAmplitude_comm,
        TNLean.PEPS.fullMultiplicityBondMap_mulVec,
        TNLean.PEPS.fullMultiplicityBondMap_apply,
        TNLean.PEPS.physicalProductMap_fullMultiplicityBondMap_apply}
    \leanok
    \uses{def:peps_parent_support_copy_amplitude,
        thm:peps_full_multiplicity_bond_map}
    Let $F=E_YJE_X^\dagger:\C^{X\times X}\to\C^{Y\times Y}$
    be the full endpoint multiplicity map, where
    $J\ket{i,a,b}=m_i^{-1/2}\sum_{u\in M_i}\ket{i,(a,u),(b,u)}$
    and an empty sum is zero. Its coefficient and action formulas are
    \begin{align}
        F_{(y,z),(x,r)} & =a(y,z)\delta_{(b(y),b(z)),(x,r)},
        \label{eq:peps_parent_support_sparse_coefficient}                        \\
        (F\phi)(y,z)    & =a(y,z)\phi(b(y),b(z)),                                \\
        (F^{\otimes E}\psi)(\beta)
                        & =\left(\prod_{f\in E}a(\beta_{h,f},\beta_{t,f})\right)
        \psi\bigl((b(\beta_{h,f}),b(\beta_{t,f}))_{f\in E}\bigr).
        \label{eq:peps_parent_support_sparse_product}
    \end{align}
    Also $a(y,z)=a(z,y)$. These formulas require neither nonempty
    multiplicity sets nor support conditions on the input vectors.
    More generally, for a coordinate inclusion $j:B\hookrightarrow A$
    with $A$ finite, its inclusion matrix satisfies
    $(E_j^\dagger\phi)(i)=\phi(j(i))$.
\end{theorem}
\begin{proof}\leanok
    The adjoint of a coordinate inclusion selects its included
    coordinates. For equal sectors, the multiplicity map retains equal
    copy labels and multiplies by $m_i^{-1/2}$; for unequal sectors its
    coefficient is zero. This proves the formula for $F\phi$.
    Evaluating it on coordinate vectors gives
    (\ref{eq:peps_parent_support_sparse_coefficient}); the product of
    its coordinate indicators selects the single input configuration
    in~(\ref{eq:peps_parent_support_sparse_product}). Equal copy labels
    have equal sectors and hence equal normalizations, proving symmetry.
\end{proof}

\begin{theorem}[Supported inverses of the multiplicity product]%
    \label{thm:peps_parent_support_retractions}
    \lean{TNLean.PEPS.fullMultiplicityBondMap_mul_adjoint_mul,
        TNLean.PEPS.physicalProductMap_fullMultiplicityBondMap_leftInverse,
        TNLean.PEPS.physicalProductMap_fullMultiplicityBondMap_rightInverse}
    \leanok
    \uses{thm:peps_full_multiplicity_bond_map}
    Assume every $M_i$ is nonempty. Let $E_X$ be the matching-sector
    inclusion and let $\mathcal F=F^{\otimes E}$ for a finite set $E$.
    Then $FF^\dagger F=F$. For $\psi\in\operatorname{im}E_X^{\otimes E}$
    and $\phi\in\operatorname{im}\mathcal F$,
    \begin{align}
        \mathcal F^\dagger\mathcal F\psi & =\psi, \\
        \mathcal F\mathcal F^\dagger\phi & =\phi.
        \label{eq:peps_parent_support_retractions}
    \end{align}
\end{theorem}
\begin{proof}\leanok
    \uses{thm:peps_full_multiplicity_bond_map,
        thm:peps_physical_product_supported_isometry}
    Since $F^\dagger F=E_XE_X^\dagger$ and $E_X^\dagger E_X=I$,
    the factorization $F=E_YJE_X^\dagger$ gives $FF^\dagger F=F$.
    Products respect multiplication. Write
    $\psi=E_X^{\otimes E}\chi$ and $\phi=\mathcal F\eta$, respectively,
    and cancel the corresponding local factors.
\end{proof}

\subsection{One-ended boundary transformations}

\begin{definition}[The boundary filter at a fixed exterior endpoint]%
    \label{def:peps_parent_support_boundary_filter}
    \lean{TNLean.PEPS.multiplicityBoundaryMap}
    \leanok
    \uses{def:peps_parent_support_copy_amplitude}
    For weights $\lambda_i\in\C$ and a fixed exterior endpoint
    $z\in Y$, define $B_z^\lambda:\C^X\to\C^Y$. For $y=(i,a,u)$, put
    \begin{align}
        (B_z^\lambda)_{y,x}
        =\lambda_i a(y,z)\delta_{b(y),x}.
        \label{eq:peps_parent_support_boundary_filter}
    \end{align}
    Write $B_z=B_z^1$ for unit weights.
\end{definition}

\begin{theorem}[Boundary filters and block matrices]%
    \label{thm:peps_parent_support_boundary_factors}
    \lean{TNLean.PEPS.multiplicityBoundaryMap_tail,
        TNLean.PEPS.multiplicityBoundaryMap_head,
        TNLean.PEPS.multiplicityBoundaryMap_mulVec,
        TNLean.PEPS.multiplicityBoundaryMap_one_conjTranspose,
        TNLean.PEPS.blockDiagonal_weight_column,
        TNLean.PEPS.blockDiagonal_weight_row}
    \leanok
    \uses{def:peps_parent_support_boundary_filter}
    For arbitrary matrices $T_i$ on $X_i$, write
    $M=\bigoplus_iT_i$, $M^{[m]}=\bigoplus_i(T_i\otimes I_{M_i})$,
    and $M_\lambda=\bigoplus_i\lambda_iT_i$.
    For $x\in X$ and $y,z\in Y$,
    \begin{align}
        \sum_{r\in Y}(B_z^\lambda)_{r,x}M^{[m]}_{r,y}
         & =a(y,z)(M_\lambda)_{x,b(y)},                     \\
        \sum_{r\in Y}M^{[m]}_{y,r}(B_z^\lambda)_{r,x}
         & =a(y,z)(M_\lambda)_{b(y),x},
        \label{eq:peps_parent_support_boundary_contraction} \\
        (B_z^\lambda\phi)(y)
         & =\lambda_{i(y)}a(y,z)\phi(b(y)).
    \end{align}
    The entries of $B_z$ are real, hence $B_z^\dagger=B_z^{\mathsf T}$.
    A block scalar may be read from either index:
    $(M_\lambda)_{x,r}=\lambda_{i(x)}M_{x,r}
        =\lambda_{i(r)}M_{x,r}$.
\end{theorem}
\begin{proof}\leanok
    The block-diagonal and identity factors force the sector and copy
    coordinates in the sum. The remaining matrix entry is that of
    $T_i$, multiplied by $\lambda_i m_i^{-1/2}$. Transposing the
    block matrices exchanges the two contraction formulas.
    The action formula follows by summing the single surviving input
    coordinate in~(\ref{eq:peps_parent_support_boundary_filter}).
    Its unit-weight coefficient is real. Off-block entries vanish;
    on a block, the row and column sector labels agree.
\end{proof}

\begin{theorem}[Adjoint boundary and internal factors]%
    \label{thm:peps_parent_support_boundary_adjoints}
    \lean{TNLean.PEPS.multiplicityBoundaryMap_adjoint_head,
        TNLean.PEPS.multiplicityBoundaryMap_adjoint_tail,
        TNLean.PEPS.fullMultiplicityBondMap_adjoint_mulVec_repeated,
        TNLean.PEPS.blockFourthRootWeight_mul_blockMatrixRepresentation,
        TNLean.PEPS.blockMatrixRepresentation_mul_blockFourthRootWeight}
    \leanok
    \uses{thm:peps_parent_support_boundary_factors}
    Suppose $\lambda_i\ne0$ for every $i$. With the preceding notation,
    \begin{align}
        \left[B_z^\dagger(M^{[m]}_{\cdot,y})\right](x)
         & =\lambda_{i(y)}^{-1}a(y,z)(M_\lambda)_{x,b(y)}, \\
        \left[B_z^\dagger(M^{[m]}_{y,\cdot})\right](x)
         & =\lambda_{i(y)}^{-1}a(y,z)(M_\lambda)_{b(y),x}.
        \label{eq:peps_parent_support_boundary_adjoint}
    \end{align}
    Independently of these weights, if every $M_i$ is nonempty then
    \begin{align}
        F^\dagger\operatorname{vec}\left(\bigoplus_i T_i\otimes I_{M_i}\right)
        =\operatorname{vec}\left(\bigoplus_i\sqrt{m_i}\,T_i\right)
        \label{eq:peps_parent_support_internal_adjoint}
    \end{align}
    for arbitrary matrices $T_i$ on $X_i$, where
    $\operatorname{vec}(T)(x,r)=T_{x,r}$.
    In particular, for representations $D_i$ of dimensions $d_i$,
    $D=\bigoplus_iD_i$ and $W=\bigoplus_i d_i^{1/4}I_{d_i}$ satisfy
    $WD(g)=D(g)W=\bigoplus_i d_i^{1/4}D_i(g)$.
\end{theorem}
\begin{proof}\leanok
    \uses{thm:peps_parent_support_boundary_factors,
        thm:peps_full_multiplicity_bond_map,
        lem:peps_matching_sector_bond_inclusion}
    Apply~(\ref{eq:peps_parent_support_boundary_contraction}) with
    unit weights and use $B_z^\dagger=B_z^{\mathsf T}$. Reading the
    block weight from the fixed $y$ index permits cancellation of
    $\lambda_{i(y)}$. The repeated matrix-entry vector in
    (\ref{eq:peps_parent_support_internal_adjoint}) is $F$ applied
    to the square-root-weighted block vector. Apply
    $F^\dagger F=E_XE_X^\dagger$;
    the block-diagonal input lies in the matching-sector range.
    The fourth-root identities follow by multiplying block-scalar
    matrices on the left and on the right.
\end{proof}

\subsection{Physical coordinates and regional operator blocks}

\begin{definition}[Whole-graph and regional bond coordinates]%
    \label{def:peps_parent_support_bond_coordinates}
    \lean{TNLean.PEPS.graphBondConfigEquiv,
        TNLean.PEPS.graphBondPhysicalEquiv,
        TNLean.PEPS.graphBondPhysicalEquiv_apply,
        TNLean.PEPS.regionBondConfigEquiv,
        TNLean.PEPS.regionBondConfigEquiv_symm_apply,
        TNLean.PEPS.regionBondPhysicalEquiv,
        TNLean.PEPS.regionOutsideBoundaryEndpoint}
    \leanok
    \uses{def:peps_region_half_edge_label_equiv}
    Choose bijections $e_v:X^{\operatorname{Inc}(v)}\to\{0,\ldots,p-1\}$.
    Decoding a numbered configuration and grouping its endpoint labels
    gives a bijection $j_e:\{0,\ldots,p-1\}^V\to(X\times X)^E$.
    The pair on $f=(t,h)$ is $(x_{h,f},x_{t,f})$.
    Its coefficient-space isomorphism is
    $J_e\psi(\beta)=\psi(j_e^{-1}\beta)$.
    Regionally the same regrouping gives
    \begin{align}
        j_{e,R}:\{0,\ldots,p-1\}^R
         & \longrightarrow X^{\partial E_R}\times(X\times X)^{E_R},
         & J_{e,R}\phi                                              & =\phi\circ j_{e,R}^{-1}.
        \label{eq:peps_parent_support_bond_coordinates}
    \end{align}
    The inverse assigns each crossing label to its inside incidence,
    and each internal pair to its head and tail incidences, before
    applying the corresponding $e_v$.
    If $\sigma$ is a numbered exterior configuration, let
    $\gamma^e_R(\sigma)_f$ be its endpoint label outside $R$:
    the head label when $t\in R$, and the tail label when $h\in R$.
\end{definition}

\begin{definition}[The global bond operator in physical coordinates]%
    \label{def:peps_parent_support_numbered_operator}
    \lean{TNLean.PEPS.graphNumberedBondMatrix}
    \leanok
    \uses{def:peps_parent_support_bond_coordinates}
    Let $e^X,e^Y$ enumerate incident configurations of $X,Y$,
    respectively, and let $F:\C^{X\times X}\to\C^{Y\times Y}$.
    Define the matrix $N_F$ between the numbered physical spaces by
    \begin{align}
        (N_F)_{\tau,\sigma}
        =\prod_{f\in E}F_{(j_{e^Y}\tau)_f,(j_{e^X}\sigma)_f}.
        \label{eq:peps_parent_support_numbered_operator}
    \end{align}
\end{definition}

\begin{theorem}[The bond operator under coordinate regrouping]%
    \label{thm:peps_parent_support_global_coordinates}
    \lean{TNLean.PEPS.graphNumberedBondMatrix_eq_submatrix,
        TNLean.PEPS.graphBondPhysicalEquiv_mulVec,
        TNLean.PEPS.graphNumberedBondMatrix_conjTranspose}
    \leanok
    \uses{def:peps_parent_support_numbered_operator}
    The matrix $N_F$ is the matrix of $F^{\otimes E}$ with rows
    relabelled by $j_{e^Y}$ and columns by $j_{e^X}$. Thus
    $J_{e^Y}N_F=F^{\otimes E}J_{e^X}$ and $N_F^\dagger=N_{F^\dagger}$,
    where the latter operator exchanges the two enumerations.
\end{theorem}
\begin{proof}\leanok
    The product coefficients agree by their definitions.
    Reindex the matrix-vector sum by the bijection $j_{e^X}$ to
    obtain the intertwining identity. Conjugating each factor in
    (\ref{eq:peps_parent_support_numbered_operator}) proves the
    adjoint identity.
\end{proof}

\begin{definition}[Blocks with fixed exterior configurations]%
    \label{def:peps_parent_support_operator_block}
    \lean{TNLean.PEPS.regionOperatorBlock}
    \leanok
    For a matrix $M$ between global numbered physical spaces, define
    its regional block by
    $M_R[\tau,\sigma]_{y,x}=M_{y\cup\tau,x\cup\sigma}$,
    where $x,y$ are configurations on $R$, and $\sigma,\tau$ are
    configurations on $R^c$. Write
    $S_{R,\sigma}\psi(x)=\psi(x\cup\sigma)$ for an exterior slice.
\end{definition}

\begin{theorem}[Regional blocks imply parent-space inclusions]%
    \label{thm:peps_parent_support_block_criterion}
    \lean{TNLean.PEPS.regionSliceMap_mulVec,
        TNLean.PEPS.regionOperatorBlock_conjTranspose,
        TNLean.PEPS.mulVec_mem_regionParentGroundSpace_of_blocks,
        TNLean.PEPS.map_regionParentGroundSpace_le_of_blocks}
    \leanok
    \uses{def:peps_parent_support_operator_block,def:peps_region_ground_space}
    For every exterior output configuration $\tau$,
    \begin{align}
        S_{R,\tau}(M\psi)
                                   & =\sum_\sigma M_R[\tau,\sigma]S_{R,\sigma}\psi,
        \label{eq:peps_parent_support_block_slice}                                  \\
        (M^\dagger)_R[\sigma,\tau] & =M_R[\tau,\sigma]^\dagger.
    \end{align}
    Let $A,B$ be tensors on the same graph, possibly with different
    physical dimensions, and let $\mathcal R=(R_i)_{i\in\iota}$ be
    any region family. If
    $M_{R_i}[\tau,\sigma]\mathcal G_{R_i}(A)
        \subseteq\mathcal G_{R_i}(B)$ for every $i,\tau,\sigma$, then
    $M\mathcal P_{\mathcal R}(A)\subseteq\mathcal P_{\mathcal R}(B)$.
    Here $\mathcal P_{\mathcal R}(A)$ consists of vectors whose every
    exterior slice on $R_i$ belongs to $\mathcal G_{R_i}(A)$.
\end{theorem}
\begin{proof}\leanok
    Split the global input-coordinate sum into its regional and
    exterior coordinates to obtain
    (\ref{eq:peps_parent_support_block_slice}). Taking the adjoint
    exchanges input and output configurations.
    If $\psi\in\mathcal P_{\mathcal R}(A)$, every summand in
    (\ref{eq:peps_parent_support_block_slice}) lies in
    $\mathcal G_{R_i}(B)$ by the assumed block inclusion. This space
    is closed under sums, proving all target parent constraints.
\end{proof}

\begin{definition}[Crossing and exterior factors of a bond operator]%
    \label{def:peps_parent_support_exterior_factors}
    \lean{TNLean.PEPS.graphBoundaryBondMatrix,
        TNLean.PEPS.graphExteriorBondCoefficient,
        TNLean.PEPS.graphMultiplicityBlockCoefficient}
    \leanok
    \uses{def:peps_parent_support_numbered_operator,
        def:peps_parent_support_copy_amplitude}
    Fix exterior configurations $\tau,\sigma$, and put
    $\gamma_Y=\gamma_R^{e^Y}(\tau)$,
    $\gamma_X=\gamma_R^{e^X}(\sigma)$.
    The matrix on the inside endpoint of a crossing edge is
    \begin{align}
        C_f(y,x)=
        \begin{cases}
            F_{(\gamma_Y(f),y),(\gamma_X(f),x)}, & t(f)\in R, \\
            F_{(y,\gamma_Y(f)),(x,\gamma_X(f))}, & h(f)\in R.
        \end{cases}
        \label{eq:peps_parent_support_crossing_matrix}
    \end{align}
    Decode the exterior endpoint labels as $\tau_{h,f},\tau_{t,f}$
    and $\sigma_{h,f},\sigma_{t,f}$. Define the exterior scalar and,
    for the multiplicity map, its crossing-indicator factor by
    \begin{align}
        c_{R,F}(\tau,\sigma)      & =\prod_{f\in E_{R^c}}
        F_{(\tau_{h,f},\tau_{t,f}),(\sigma_{h,f},\sigma_{t,f})}, \\
        \widehat c_R(\tau,\sigma) & =c_{R,F}(\tau,\sigma)
        \prod_{f\in\partial E_R}\delta_{b(\gamma_Y(f)),\gamma_X(f)}.
    \end{align}
\end{definition}

\begin{theorem}[Factorization of regional bond-operator blocks]%
    \label{thm:peps_parent_support_bond_blocks}
    \lean{TNLean.PEPS.prod_edge_eq_exterior_boundary_internal,
        TNLean.PEPS.regionOperatorBlock_graphNumberedBondMatrix}
    \leanok
    \uses{def:peps_parent_support_exterior_factors,
        def:peps_parent_support_mixed_product,
        def:peps_parent_support_operator_block}
    For any weights $w_f$, the edge partition gives
    $\prod_{f\in E}w_f=(\prod_{f\in E_{R^c}}w_f)
        (\prod_{f\in\partial E_R}w_f)(\prod_{f\in E_R}w_f)$.
    Accordingly, the blocks of the global bond operator satisfy
    \begin{align}
        (N_F)_R[\tau,\sigma]_{y,x}
        =c_{R,F}(\tau,\sigma)\cdot
        \mathcal M((C_f)_{f\in\partial E_R},(F)_{f\in E_R})_
        {j_{e^Y,R}(y),j_{e^X,R}(x)}.
        \label{eq:peps_parent_support_bond_blocks}
    \end{align}
\end{theorem}
\begin{proof}\leanok
    Every edge has zero, one, or two endpoints in $R$. Separate the
    product in~(\ref{eq:peps_parent_support_numbered_operator}) into
    these three cases. The first contributes $c_{R,F}$, the second
    gives~(\ref{eq:peps_parent_support_crossing_matrix}), and the third
    retains $F$ on its head--tail pair. Regional regrouping places
    these factors in the two coordinate families of $\mathcal M$.
\end{proof}

\begin{theorem}[Regional blocks of the multiplicity map]%
    \label{thm:peps_parent_support_copy_blocks}
    \lean{TNLean.PEPS.fullMultiplicityBondMap_fixed_tail,
        TNLean.PEPS.fullMultiplicityBondMap_fixed_head,
        TNLean.PEPS.graphBoundaryBondMatrix_fullMultiplicityBondMap,
        TNLean.PEPS.regionOperatorBlock_fullMultiplicityBondMap}
    \leanok
    \uses{def:peps_parent_support_boundary_filter,
        def:peps_parent_support_exterior_factors}
    For the multiplicity map $F$ and any $z,y\in Y$, $r,x\in X$,
    \begin{align}
        F_{(y,z),(x,r)} & =\delta_{b(z),r}(B_z)_{y,x}, \\
        F_{(z,y),(r,x)} & =\delta_{b(z),r}(B_z)_{y,x}.
        \label{eq:peps_parent_support_fixed_endpoint}
    \end{align}
    Thus $C_f=\delta_{b(\gamma_Y(f)),\gamma_X(f)}B_{\gamma_Y(f)}$
    in either orientation, and
    \begin{align}
        (N_F)_R[\tau,\sigma]_{y,x}
        =\widehat c_R(\tau,\sigma)\cdot
        \mathcal M((B_{\gamma_Y(f)})_{f\in\partial E_R},(F)_{f\in E_R})_
        {j_{e^Y,R}(y),j_{e^X,R}(x)}.
        \label{eq:peps_parent_support_copy_blocks}
    \end{align}
\end{theorem}
\begin{proof}\leanok
    \uses{thm:peps_parent_support_sparse_action,
        thm:peps_parent_support_bond_blocks}
    In~(\ref{eq:peps_parent_support_sparse_coefficient}), separate
    the indicator for the fixed endpoint from the inside endpoint.
    Symmetry of $a(y,z)$ gives the other orientation.
    Substitute~(\ref{eq:peps_parent_support_fixed_endpoint}) into
    (\ref{eq:peps_parent_support_bond_blocks}) and collect the
    crossing indicators into $\widehat c_R$.
\end{proof}

\subsection{Actual open contractions and their product support}

\begin{definition}[Boundary and internal factors of an open contraction]%
    \label{def:peps_parent_support_open_factors}
    \lean{TNLean.PEPS.graphOpenBoundaryFactor,
        TNLean.PEPS.graphOpenInternalFactor,
        TNLean.PEPS.graphOpenBondCoordinates,
        TNLean.PEPS.graphOpenBondCoordinates_eq_sum_prod}
    \leanok
    Let $G$ and $X$ be finite, $U:G\to\operatorname{Mat}_X(\C)$ a
    representation, and $W$ a matrix on $X$. For $q:R\to G$, define
    \begin{align}
        b_f(q;\theta,x) & =
        \begin{cases}
            (WU(q_t^{-1}))_{\theta,x}, & t\in R, \\
            (U(q_h)W)_{x,\theta},      & h\in R,
        \end{cases}              \\
        k_f(q;x_h,x_t)  & =(W^2U(q_hq_t^{-1}))_{x_h,x_t}.
    \end{align}
    With virtual boundary labels $\theta\in X^{\partial E_R}$, put
    \begin{align}
        H_R^{U,W}(\theta)(\sigma,\beta)
        =\sum_{q:R\to G}|G|^{-|R|}
        \prod_{f\in\partial E_R}b_f(q;\theta_f,\sigma_f)
        \prod_{f\in E_R}k_f(q;\beta_{h,f},\beta_{t,f}).
        \label{eq:peps_parent_support_open_factors}
    \end{align}
    This is the coherent open bond vector in crossing and internal
    head--tail coordinates.
\end{definition}

\begin{definition}[Numbered canonical tensors and open bond spaces]%
    \label{def:peps_parent_support_canonical_tensor}
    \lean{TNLean.PEPS.numberedGraphDressedAveragingTensor,
        TNLean.PEPS.graphOpenBondSpace,
        TNLean.PEPS.graphDressedAveragingSite_one}
    \leanok
    \uses{def:peps_parent_support_bond_coordinates,
        def:peps_parent_support_open_factors}
    Let $a_v^{U,W}(\eta,\xi)$ be the dressed averaging site. In a
    chosen incident enumeration $e$, the numbered tensor is
    $A(U,W,e)_v(\eta,s)=a_v^{U,W}(\eta,e_v^{-1}(s))$, with virtual
    labels numbered by a bijection $X\cong\{0,\ldots,|X|-1\}$.
    Dressing by $W=I$ gives the unweighted averaging site.
    Define the open bond space by
    $\mathcal H_R(U,W)=\operatorname{span}_{\C}
        \{H_R^{U,W}(\theta):\theta\in X^{\partial E_R}\}$.
\end{definition}

\begin{theorem}[Numbered regional ranges in open bond coordinates]%
    \label{thm:peps_parent_support_open_coordinates}
    \lean{TNLean.PEPS.regionBondPhysicalEquiv_openRegionWeight,
        TNLean.PEPS.map_regionGroundSpace_eq_graphOpenBondSpace,
        TNLean.PEPS.regionBondPhysicalEquiv_mem_graphOpenBondSpace_iff,
        TNLean.PEPS.regionInternalBondSlice_openRegionWeight}
    \leanok
    \uses{def:peps_parent_support_canonical_tensor,
        def:peps_region_ground_space,def:peps_open_sector_slice}
    Suppose $WU(g)=U(g)W$ for every $g$. Then
    \begin{align}
        J_{e,R}\mathcal C_R(A(U,W,e))\ket{\theta}
                                      & =H_R^{U,W}(\theta), \\
        J_{e,R}\mathcal G_R(A(U,W,e)) & =\mathcal H_R(U,W).
        \label{eq:peps_parent_support_open_coordinates}
    \end{align}
    In particular, $J_{e,R}\phi\in\mathcal H_R(U,W)$ if and only if
    $\phi\in\mathcal G_R(A(U,W,e))$.
    More generally, for any site family $a_v(\eta,\xi)$, numbering
    its virtual and physical labels leaves every internal slice of
    each open column unchanged: with the same $\theta$ and fixed
    crossing physical labels $\sigma$, the numbered and incidence
    contractions give the same function of internal head--tail pairs.
\end{theorem}
\begin{proof}\leanok
    \uses{thm:peps_graph_open_averaging_bond_state}
    Decode the numbered labels and apply the actual open contraction
    formula. The inverse regional coordinate map puts each crossing,
    head, and tail label back at its original incidence, proving the
    column equality in~(\ref{eq:peps_parent_support_open_coordinates}).
    Open columns span $\mathcal G_R$, and their boundary numberings
    are bijective, so their images span exactly $\mathcal H_R$.
    Injectivity of $J_{e,R}$ gives the membership equivalence.
    The final assertion follows directly by decoding the virtual
    boundary labels in the finite contraction sum and then fixing
    the crossing physical coordinates; it needs no representation
    or commutation hypothesis.
\end{proof}

\begin{theorem}[Edge coverage promotes internal support to global support]%
    \label{thm:peps_parent_support_internal_slices}
    \lean{TNLean.PEPS.regionInternalBondSlice_regionSliceMap_update,
        TNLean.PEPS.regionParentGroundSpace_mem_product_range_of_internal_slices}
    \leanok
    \uses{def:peps_parent_support_bond_coordinates,
        def:peps_open_sector_slice}
    Let $A$ be any tensor with incident enumeration $e$ by a finite
    endpoint set $X$, and let $F:\C^Z\to\C^{X\times X}$ with $Z$
    finite. Suppose every edge is internal to some region $R_i$.
    Assume that, for every $i$, every crossing physical configuration
    $\sigma$, and every $\phi\in\mathcal G_{R_i}(A)$,
    \begin{align}
        \mathcal S_{R_i,\sigma}\phi
        \in\operatorname{im}F^{\otimes E_{R_i}}.
        \label{eq:peps_parent_support_internal_hypothesis}
    \end{align}
    Then $J_e\psi\in\operatorname{im}F^{\otimes E}$ for every
    $\psi\in\mathcal P_{\mathcal R}(A)$.
    The coordinate identity used here is the following. Fix global
    bond labels $\beta$ and an internal edge $f\in E_R$. Let
    $\sigma$ be their crossing labels inside $R$, and let $\tau$
    be their numbered exterior configuration. For every endpoint pair $s$,
    \begin{align}
        (\mathcal S_{R,\sigma}S_{R,\tau}\psi)
        (\beta|_{E_R}[f\mapsto s])
        =(J_e\psi)(\beta[f\mapsto s]).
        \label{eq:peps_parent_support_slice_update}
    \end{align}
\end{theorem}
\begin{proof}\leanok
    \uses{thm:peps_parent_support_product_range}
    Updating an internal edge changes only its two inside incidences;
    it leaves all crossing and exterior incidences fixed. This proves
    (\ref{eq:peps_parent_support_slice_update}). For each edge $f$,
    choose a covering region $R_i$. The parent condition places
    $S_{R_i,\tau}\psi$ in $\mathcal G_{R_i}(A)$, so
    (\ref{eq:peps_parent_support_internal_hypothesis}) and the forward
    implication of~(\ref{eq:peps_parent_support_product_range}) put its
    $f$-slice in $\operatorname{im}F$. By
    (\ref{eq:peps_parent_support_slice_update}) this is the global
    $f$-slice. Apply the reverse implication of
    (\ref{eq:peps_parent_support_product_range}) to all edges.
\end{proof}

\begin{theorem}[Bond-factor support of canonical parent vectors]%
    \label{thm:peps_parent_support_factor_ranges}
    \lean{TNLean.PEPS.graphRegionInternalBondSlice_mem_range_of_factors,
        TNLean.PEPS.regionInternalBondSlice_mem_range_of_factors,
        TNLean.PEPS.regionParentGroundSpace_mem_product_range_of_factors}
    \leanok
    \uses{def:peps_parent_support_canonical_tensor,
        def:peps_open_sector_slice}
    Let $U:G\to\operatorname{Mat}_X(\C)$ be a representation of a
    finite group, let $W$ commute with every $U(g)$, and let
    $F:\C^Z\to\C^{X\times X}$, with $X,Z$ finite. Suppose
    $\operatorname{vec}(W^2U(g))\in\operatorname{im}F$ for every $g$.
    Every internal slice of every open column, with arbitrary virtual
    boundary labels, belongs to $\operatorname{im}F^{\otimes E_R}$.
    The same holds for every $\phi\in\mathcal G_R(A(U,W,e))$.
    If every edge is internal to a member of $\mathcal R$, then
    \begin{align}
        J_e\mathcal P_{\mathcal R}(A(U,W,e))
        \subseteq\operatorname{im}F^{\otimes E}.
        \label{eq:peps_parent_support_factor_ranges}
    \end{align}
\end{theorem}
\begin{proof}\leanok
    \uses{thm:peps_open_sector_weighted_slice,
        thm:peps_parent_support_internal_slices,
        thm:peps_parent_support_open_coordinates}
    Choose $x_g$ with $Fx_g=\operatorname{vec}(W^2U(g))$.
    For each internal slice use the preimage
    \begin{align}
        \chi(\alpha)=\sum_{q:R\to G}c_R(\theta,\sigma;q)
        \prod_{f=(t,h)\in E_R}x_{q_hq_t^{-1}}(\alpha_f),
    \end{align}
    where $c_R$ retains all boundary factors and the averaging
    normalization. Applying $F^{\otimes E_R}$ term by term gives
    the weighted slice formula. Numbered and incidence slices agree,
    and the regional range is spanned by open columns, so the
    inclusion extends by linearity. Theorem~\ref{thm:peps_parent_support_internal_slices}
    then proves~(\ref{eq:peps_parent_support_factor_ranges}).
\end{proof}

\begin{theorem}[Product multiplicity support of the repeated parent]%
    \label{thm:peps_parent_support_repeated_ranges}
    \lean{TNLean.PEPS.graphRegionInternalBondSlice_multiplicityRestored_mem_range,
        TNLean.PEPS.regionInternalBondSlice_mem_fullMultiplicityRange,
        TNLean.PEPS.regionParentGroundSpace_mem_product_fullMultiplicityRange}
    \leanok
    \uses{def:peps_parent_support_canonical_tensor,
        thm:peps_full_multiplicity_bond_map}
    Let $d_i>0$ and let $D_i:G\to\operatorname{Mat}_{d_i}(\C)$ be
    representations of a finite group. Put
    $U'(g)=\bigoplus_i(D_i(g)\otimes I_{d_i})$, and let $F$ restore
    multiplicity $d_i$ on the full endpoint pairs. Every internal
    slice of each actual open column for the unweighted averaging
    tensor of $U'$ lies in $\operatorname{im}F^{\otimes E_R}$.
    This also holds for each vector in the numbered regional range
    $\mathcal G_R(A(U',I,e))$.
    For any edge-covering family $\mathcal R$,
    \begin{align}
        J_e\mathcal P_{\mathcal R}(A(U',I,e))
        \subseteq\operatorname{im}F^{\otimes E}.
        \label{eq:peps_parent_support_repeated_ranges}
    \end{align}
    Neither irreducibility nor unitarity of the $D_i$ is assumed.
\end{theorem}
\begin{proof}\leanok
    \uses{thm:peps_open_sector_weighted_slice,
        thm:peps_coherent_multiplicity_transport,
        thm:peps_parent_support_open_coordinates,
        thm:peps_parent_support_internal_slices}
    Expand each target slice with its own boundary coefficients
    $c_R^{U',I}(\theta,\sigma;q)$. Coherent multiplicity restoration
    sends the internal factors
    $\operatorname{vec}(\bigoplus_i\sqrt{d_i}\,D_i(q_hq_t^{-1}))$
    to $\operatorname{vec}(U'(q_hq_t^{-1}))$, while retaining these
    same coefficients. This supplies a preimage for each target
    open slice. Linear combinations give all regional vectors;
    edge coverage and
    Theorem~\ref{thm:peps_parent_support_internal_slices} give
    (\ref{eq:peps_parent_support_repeated_ranges}).
\end{proof}
